\documentclass[12pt,a4paper]{amsart}

\usepackage{amsmath,amssymb,amsthm}
\usepackage{mathtools}
\usepackage[colorlinks=true,linkcolor=blue!60!black,
citecolor=green!50!black,urlcolor=blue!70!black]{hyperref}
\usepackage[shortlabels]{enumitem}
\usepackage{booktabs}

%% ── Theorem environments ─────────────────────────────────────────
\newtheorem{theorem}{Theorem}[section]
\newtheorem{lemma}[theorem]{Lemma}
\newtheorem{proposition}[theorem]{Proposition}
\newtheorem{corollary}[theorem]{Corollary}
\theoremstyle{definition}
\newtheorem{definition}[theorem]{Definition}
\newtheorem{remark}[theorem]{Remark}
\newtheorem{problem}{Problem}[section]
\newtheorem{assumption}[theorem]{Assumption}

%% ── Notation (aligned with Paper I) ──────────────────────────────
\newcommand{\PW}{\mathrm{PW}_\omega}
\newcommand{\Kop}{\mathcal{K}_c}
\newcommand{\GVM}{G^{(N)}_{\mathbf{p}}}
\newcommand{\GVMk}{G^{(N_k)}_{\mathbf{p}_k}}
\newcommand{\GVMkw}{\widetilde{G}^{(N_k)}_{\mathbf{p}_k}}
\newcommand{\psin}{\psi_n^{(c)}}
\newcommand{\lam}{\lambda_n(c)}
\newcommand{\norm}[1]{\|#1\|}
\newcommand{\ip}[2]{\langle #1,\, #2\rangle}
\newcommand{\EN}{E_N}
\newcommand{\Lcal}{\mathcal{L}^1_{\mathrm{sa}}}
\newcommand{\Gcal}{\widehat{\mathcal{G}}}

%% ── Paper metadata ───────────────────────────────────────────────
\title[Scaling Limits, Uniform Coercivity, and Trace Formula]
{Scaling Limits of Prolate Gram Forms,\\
Uniform Coercivity, and a Trace Formula\\
Toward Weil Positivity\\[4pt]
\large (Paper~II in the Prolate--Weil Program)}

\author{Anonymous Author}
\address{}
\email{}
\date{\today}

\subjclass[2020]{%
  42C05,
  42B10,
  47B10,
  47B32,
  11M06,
  11N05
}

\keywords{prolate spheroidal wave functions, Gram form, scaling limit,
uniform coercivity, trace formula, Weil positivity, prime number theorem,
Connes--Consani--Moscovici, trace-class operators}

\begin{document}

\begin{abstract}
This paper continues the program initiated in \cite{PaperI}, where the
prolate Gram form $Q^{(N)}$ was shown to be unconditionally coercive on
the $N$-dimensional prolate subspace $V_N$, with explicit spectral
stability under the Discrete Spectral Transfer Property (DSTP)
of \cite{PaperI}.
In the current Reinschrift of Paper~I, the DSTP is verified
for random sampling and Gauss--PSWF nodes, while prime sampling is left open.

For consistency with \cite{PaperI}, all references to the spectral decay
condition formerly labelled~(T.2) in earlier drafts are replaced by the
DSTP throughout this paper.

Here we study the coupled scaling limit
\[
  N_k \to \infty,\quad c_k \to \infty,\quad
  \frac{N_k}{c_k} \to \rho \in \Bigl(0,\,\frac{2}{\pi}\Bigr),
\]
for sequences of prime sampling sets, under an explicit conditional
hypothesis~\textup{(H)} asserting that the prime sets satisfy a quantitative
version of the DSTP along the sequence.

We establish three main results.
\begin{enumerate}[(I)]
\item \textbf{(Scaling limit, \S\ref{sec:scaling}.)} We define a precise notion of
$\rho$-admissible scaling sequence, isolate the conditional prime-sampling
input in Assumption~\ref{ass:primeT2}, and identify the natural operator
space for convergence as $\Lcal(L^2(\mathbb{R}))$ with the weak-$*$ topology.
\item \textbf{(Compactness and conditional coercivity, \S\ref{sec:coercivity}.)}
Compactness of normalized Gram operators is unconditional once a
$\rho$-admissible sequence is given, while the convergence
$\lambda_{\min}(\GVMk)\to1$ is proved conditionally on~\textup{(H)}.
\item \textbf{(Trace formula, \S\ref{sec:trace}.)} We establish asymptotic trace
formulae for the unweighted and weighted Gram matrices, connecting the
weighted trace to the prime number theorem. Every weak-$*$ accumulation point
$\Gcal_\infty$ is positive with $\operatorname{tr}(\Gcal_\infty)\leq1$.
\end{enumerate}
The characterization of $\Gcal_\infty$ as a Weil operator, the equality
$\operatorname{tr}(\Gcal_\infty)=1$, and the verification of Assumption~\textup{(H)}
for prime sampling remain open.
No claim is made about zeros of~$\zeta(s)$.
\end{abstract}

\maketitle
\tableofcontents

\section{Introduction}\label{sec:intro}

\subsection{Background and motivation}

The \emph{Weil positivity criterion} \cite{Weil1952,Bombieri2000}
states that the Riemann Hypothesis is equivalent to positivity of an
explicit quadratic form built from primes and archimedean terms.
The program of Connes--Consani--Moscovici \cite{CCM_PNAS2022,CCM2025}
seeks a spectral realization via the prolate operator $\Kop$.

In \cite{PaperI} the prolate Gram form
$Q^{(N)}(f)=\sum_{j=1}^N|f(p_j)|^2$ on $V_N$ was shown to be:
\begin{itemize}
\item \emph{unconditionally coercive} once the sampling set is unisolvent;
\item \emph{explicitly stable and Weil-decomposable} under the
Discrete Spectral Transfer Property (DSTP) of \cite{PaperI}.
\end{itemize}
Moreover, the present Reinschrift of Paper~I proves the relevant finite-dimensional
results for random sampling and Gauss--PSWF nodes, but does \emph{not} prove
the DSTP for prime sampling; that issue is left open there.
The purpose of the present paper is therefore to analyze the coupled limit
\(N_k,c_k\to\infty\) in a way that cleanly separates the unconditional
operator-theoretic structure from the conditional arithmetic input.

\medskip
\noindent\textbf{Core outcome.}
This paper constructs the natural normalized finite-rank Gram operators
attached to prime sampling in prolate spaces and studies their scaling limits.
The compactness statement and the existence of positive trace-class
accumulation points are proved at the operator-theoretic level, while the
strong coercivity conclusion $\lambda_{\min}(\GVMk)\to1$ is obtained only
under an explicit prime-sampling hypothesis \textup{(H)}; see
Assumption~\ref{ass:primeT2}.

\subsection{Two-level structure}

We maintain a strict distinction between unconditional and conditional results:

\begin{enumerate}[(A)]
\item \emph{Operator-theoretic layer.} Once a sequence of normalized finite-rank
Gram operators is given with uniformly bounded trace norm, compactness and the
existence of weak-$*$ accumulation points follow by general functional analysis.
The prime number theorem also yields the structural constraint $\omega_k\to0$
in the prime-window regime.
\item \emph{Arithmetic--sampling layer.} Any quantitative convergence statement for
$\lambda_{\min}(\GVMk)$ requires a spectral comparison estimate of DSTP type
along the prime sequence. Since this is open in Paper~I, we isolate it here
as an explicit assumption rather than importing it as a theorem.
\end{enumerate}

\subsection{What this paper does not claim}

No statement is made about zeros of $\zeta(s)$.
We do not prove that prime sampling satisfies the DSTP of \cite{PaperI}.
We do not prove that the limit operator $\Gcal_\infty$ is uniquely determined,
nor that $\operatorname{tr}(\Gcal_\infty)=1$.
We also do not prove a Mellin-isometric embedding identifying
$\Gcal_\infty$ with a Weil operator.

\subsection{Organization}

Section~\ref{sec:prelim} recalls the relevant inputs from Paper~I.
Section~\ref{sec:scaling} defines the admissible scaling framework and isolates
Assumption~\ref{ass:primeT2}.
Section~\ref{sec:coercivity} proves compactness and conditional uniform coercivity.
Section~\ref{sec:trace} establishes the trace formulae.
Section~\ref{sec:open} formulates the remaining open problems.

\section{Preliminaries}\label{sec:prelim}

\subsection{Notation}

We use the engineering Fourier convention
$\widehat{f}(\xi)=\int f(x)e^{-2\pi ix\xi}dx$ \cite{BonamiKaroui,PaperI}.
Bandwidth $\omega>0$, time window $T>0$, and $c=\omega T$.
The Paley--Wiener space $\PW$, the prolate concentration operator $\Kop$,
the PSWF basis $\{\psin\}$, the eigenvalues $\lambda_n(c)$, the subspace
$V_N=\mathrm{span}\{\psi_0^{(c)},\ldots,\psi_{N-1}^{(c)}\}$, and the Gram matrix
$\GVM$ are used exactly as in \cite{PaperI}.

\subsection{Input from Paper~I}

The following summary is the only part of Paper~I used in this paper.
It deliberately reflects the Reinschrift status of that paper.

\begin{theorem}[Input from Paper~I]\label{thm:paperI}
Let $N\geq1$, $c=\omega T$, and let $p_1<\cdots<p_N$ be distinct points in
$[-T,T]$. Then:
\begin{enumerate}[(I)]
\item If the sampling set is unisolvent for $V_N$, then the Gram form is coercive:
$Q^{(N)}(f)\geq \alpha_N^2\norm{f}_{L^2}^2$ for all $f\in V_N$, where
$\alpha_N=\sqrt{\lambda_{\min}(\GVM)}>0$.
\item Unconditionally, one has the continuity bound
$Q^{(N)}(f)\leq N(\omega/\pi)\norm{f}_{L^2}^2$.
\item Under the DSTP of \cite{PaperI}, one has the spectral stability bound
$\lambda_{\min}(\GVM)\geq\lambda_{N-1}(c)-\varepsilon_N$.
\item Under the DSTP of \cite{PaperI}, one has the comparison formula
$Q^{(N)}(f)=\ip{\Kop f}{f}_{L^2}+R_N(f)$ with
$|R_N(f)|\leq \varepsilon_N\norm{f}_{L^2}^2$.
\end{enumerate}
Moreover, in the Reinschrift of Paper~I, the DSTP is established for
random sampling and Gauss--PSWF nodes, while prime sampling is left open.
\end{theorem}

\begin{remark}[Landau--Widom asymptotic]\label{rem:LW}
For every fixed $\varepsilon\in(0,2/\pi)$ there exist constants
$C(\varepsilon),\eta(\varepsilon)>0$ such that
\[
  1-C(\varepsilon)e^{-\eta(\varepsilon)c}
  \le \lambda_{N-1}(c) \le 1
  \qquad\text{whenever } N\leq \bigl(\tfrac{2}{\pi}-\varepsilon\bigr)c.
\]
In particular, along any sequence with $N_k/c_k\to\rho\in(0,2/\pi)$,
one has $\lambda_{N_k-1}(c_k)\to1$.
\end{remark}

\section{The Scaling Limit}\label{sec:scaling}

\subsection{Admissible scaling sequences}

\begin{definition}[$\rho$-admissible prime sequence]\label{def:admissible}
Fix $\rho\in(0,2/\pi)$. A sequence
$(N_k,c_k,\omega_k,T_k,\mathbf{p}_k)_{k\geq1}$ is called
\emph{$\rho$-admissible for prime sampling} if:
\begin{enumerate}[(i)]
\item $N_k\in\mathbb N$ and $N_k\to\infty$;
\item $c_k=\omega_kT_k>0$ and $c_k\to\infty$;
\item $N_k/c_k\to\rho$;
\item there exists $\varepsilon_0\in(0,1)$ such that
$N_k\le (2/\pi-\varepsilon_0)c_k$ for all sufficiently large $k$;
\item $\mathbf p_k=(p_1^{(k)},\ldots,p_{N_k}^{(k)})$ consists of the first
$N_k$ prime numbers, with $p_{N_k}^{(k)}\le T_k$;
\item the associated normalized Gram operators satisfy the trace-norm boundedness
required in Lemma~\ref{lem:embed}.
\end{enumerate}
\end{definition}

\begin{remark}[Structural constraint: $\omega_k\to0$]\label{rem:omega_to_zero}
Conditions (iii) and (v) force $\omega_k\to0$.
Indeed, by the prime number theorem one has
$p_{N_k}\sim N_k\log N_k$, so condition (v) implies
$T_k\gtrsim N_k\log N_k$.
Since $c_k=\omega_kT_k\sim N_k/\rho$, we obtain
\[
  \omega_k \sim \frac{N_k}{\rho T_k}
  \lesssim \frac{1}{\rho\log N_k}\to0.
\]
Thus the coupled prime-window regime necessarily drives the bandwidth to zero;
this is a structural feature of the scaling, not a pathology.
\end{remark}

\begin{assumption}[Prime-sampling comparison hypothesis]\label{ass:primeT2}
Let $(N_k,c_k,\omega_k,T_k,\mathbf p_k)$ be a $\rho$-admissible prime sequence.
We say that Assumption~\textup{(H)} holds along the sequence if there exist numbers
$\varepsilon_{N_k}\downarrow0$ such that the prime sampling sets $\mathbf p_k$
satisfy the DSTP of \cite{PaperI} with defect level $\varepsilon_{N_k}$
for every $k$.
When a quantitative rate is needed, we additionally assume that
$\varepsilon_{N_k}=o(1)$ is sufficiently strong for the specific estimate under discussion.
\end{assumption}

\begin{remark}[Why Assumption~\textup{(H)} is explicit]
In the Reinschrift of Paper~I, prime sampling is not proved to satisfy the DSTP;
this is recorded there as an open problem.
Accordingly, any use of a prime-sampling comparison estimate in the present paper
must be conditional.
Assumption~\textup{(H)} is precisely the missing arithmetic--sampling input.
\end{remark}

\begin{proposition}[Conditional existence of admissible examples]
\label{prop:admissible_exist}
Fix $\rho\in(0,2/\pi)$. Suppose there exists a sequence of prime sampling sets for which
Assumption~\textup{(H)} holds. Then there exists a $\rho$-admissible prime sequence.
In particular, one may take $N_k=k$, choose $T_k=p_k$, set $c_k=N_k/\rho$, and define
$\omega_k=c_k/T_k$.
\end{proposition}

\begin{proof}
Set $N_k=k$, $T_k=p_k$, $c_k=N_k/\rho$, and $\omega_k=c_k/T_k$.
Then conditions (i)--(v) in Definition~\ref{def:admissible} are immediate.
The remaining trace-norm boundedness in (vi) follows once the normalized Gram operators
are defined and estimated as in Lemma~\ref{lem:embed}.
The only genuinely arithmetic comparison input is Assumption~\textup{(H)},
which is assumed rather than derived.
\end{proof}

\subsection{The natural operator space}

\begin{definition}[Embedded normalized Gram operator]\label{def:embed}
For a $\rho$-admissible prime sequence, define
\[
  \widetilde{\mathcal G}_k
  :=
  \frac{1}{\operatorname{tr}(\GVMk)}
  \sum_{m,n=0}^{N_k-1}(\GVMk)_{mn}
  |\psi_m^{(c_k)}\rangle\langle\psi_n^{(c_k)}|.
\]
This is a finite-rank positive self-adjoint operator on $L^2(\mathbb R)$
with trace exactly equal to~$1$.
\end{definition}

\begin{lemma}[Trace-class embedding]\label{lem:embed}
For each $k$:
\begin{enumerate}[(i)]
\item $\widetilde{\mathcal G}_k\ge0$.
\item $\widetilde{\mathcal G}_k\in\Lcal(L^2(\mathbb R))$.
\item $\norm{\widetilde{\mathcal G}_k}_1=1$.
\end{enumerate}
\end{lemma}

\begin{proof}
Since $\GVMk=M_k^*M_k\ge0$, the operator
$\sum_{m,n}(\GVMk)_{mn}|\psi_m^{(c_k)}\rangle\langle\psi_n^{(c_k)}|$
is positive, hence so is $\widetilde{\mathcal G}_k$.
It has rank at most $N_k$, hence it is trace class.
By construction,
\[
  \operatorname{tr}(\widetilde{\mathcal G}_k)
  =\frac{\operatorname{tr}(\GVMk)}{\operatorname{tr}(\GVMk)}=1.
\]
For positive trace-class operators, the trace norm equals the trace.
\end{proof}

\begin{theorem}[Compactness and accumulation points]\label{thm:compact}
Let $(N_k,c_k,\omega_k,T_k,\mathbf p_k)$ be a $\rho$-admissible prime sequence.
Then every subsequence of $(\widetilde{\mathcal G}_k)$ has a further subsequence converging
in the weak-$*$ topology of $\Lcal$ to some
$\Gcal_\infty\in\Lcal(L^2(\mathbb R))$ with
$\Gcal_\infty\ge0$ and $\operatorname{tr}(\Gcal_\infty)\le1$.
\end{theorem}

\begin{proof}
By Lemma~\ref{lem:embed}, the sequence lies in the positive part of the unit sphere
of $\Lcal$.
The unit ball of $\Lcal$ is weak-$*$ compact by Banach--Alaoglu.
Hence every subsequence admits a further weak-$*$ convergent subsequence.
Positivity is preserved under weak-$*$ limits, and the trace is weak-$*$
lower semicontinuous on positive trace-class operators, giving
$\operatorname{tr}(\Gcal_\infty)\le1$.
\end{proof}

\begin{remark}[Why we only obtain $\operatorname{tr}(\Gcal_\infty)\le1$]
\label{rem:trace_eq}
Weak-$*$ convergence alone does not prevent loss of trace mass at infinity.
To prove $\operatorname{tr}(\Gcal_\infty)=1$, one would need an additional tightness or
equi-integrability statement in trace norm.
That issue is genuinely open in the present framework and is recorded later as an open problem.
\end{remark}

\section{Compactness and Conditional Coercivity}\label{sec:coercivity}

\begin{theorem}[Conditional uniform coercivity]\label{thm:coercivity}
Let $(N_k,c_k,\omega_k,T_k,\mathbf p_k)$ be a $\rho$-admissible prime sequence.
Then:
\begin{enumerate}[(I)]
\item \textbf{Unconditional spectral input:}
$\lambda_{N_k-1}(c_k)\to1$.
\item \textbf{Conditional comparison estimate:}
If Assumption~\textup{(H)} holds, then
\[
  \lambda_{\min}(\GVMk)
  \ge \lambda_{N_k-1}(c_k)-\varepsilon_{N_k}
  \to1.
\]
\item \textbf{Conditional frame asymptotics:}
Under Assumption~\textup{(H)},
$\alpha_{N_k}:=\sqrt{\lambda_{\min}(\GVMk)}\to1$.
\end{enumerate}
\end{theorem}

\begin{proof}
Part (I) is exactly the Landau--Widom asymptotic from Remark~\ref{rem:LW},
using only the condition $N_k/c_k\to\rho\in(0,2/\pi)$.

For part (II), Assumption~\textup{(H)} provides the DSTP of \cite{PaperI}
along the sequence, so Theorem~\ref{thm:paperI}(III) applies at each level $k$
and yields
\[
  \lambda_{\min}(\GVMk)
  \ge \lambda_{N_k-1}(c_k)-\varepsilon_{N_k}.
\]
Since $\lambda_{N_k-1}(c_k)\to1$ and $\varepsilon_{N_k}\to0$, the conclusion follows.

Part (III) is immediate from part (II).
\end{proof}

\begin{remark}[What is unconditional and what is not]
The approach of this paper yields an unconditional compactness theorem for the
normalized operators once the scaling sequence is fixed, and an unconditional
asymptotic for the reference PSWF eigenvalues.
By contrast, any statement about the actual prime-sampled Gram matrices beyond those
formal operator properties requires Assumption~\textup{(H)}.
This is exactly where the unresolved content of prime sampling enters.
\end{remark}

\section{Trace Formula and Prime Number Consistency}\label{sec:trace}

\subsection{Trace identities}

\begin{lemma}[Dual trace identity]\label{lem:dual_trace}
For every $k$ one has:
\begin{enumerate}[(i)]
\item $\operatorname{tr}(\GVMk)=\sum_{j=1}^{N_k}K_{N_k}(p_j^{(k)},p_j^{(k)})$.
\item If Assumption~\textup{(H)} holds, then
\[
  \operatorname{tr}(\GVMk)
  =\sum_{n=0}^{N_k-1}\lambda_n(c_k)+O(N_k\varepsilon_{N_k}).
\]
\end{enumerate}
\end{lemma}

\begin{proof}
Part (i) is the standard diagonal trace computation:
\[
  \operatorname{tr}(\GVMk)
  =\sum_{n=0}^{N_k-1}(\GVMk)_{nn}
  =\sum_{j=1}^{N_k}\sum_{n=0}^{N_k-1}|\psi_n^{(c_k)}(p_j^{(k)})|^2
  =\sum_{j=1}^{N_k}K_{N_k}(p_j^{(k)},p_j^{(k)}).
\]
Part (ii) follows from writing
$\GVMk=G^{\mathrm{ref}}_{N_k,c_k}+E_{N_k}$ under the DSTP of \cite{PaperI},
taking traces, and estimating
$|\operatorname{tr}(E_{N_k})|\le N_k\norm{E_{N_k}}_2\le N_k\varepsilon_{N_k}$.
\end{proof}

\subsection{Weighted Gram matrix}

\begin{definition}[Weighted Gram matrix]\label{def:weighted}
For prime sampling, define
\[
  \GVMkw:=M_k^*\,\mathrm{diag}(\log p_1^{(k)},\ldots,\log p_{N_k}^{(k)})\,M_k.
\]
Its associated quadratic form is
\[
  \widetilde Q^{(N_k)}(f)=\sum_{j=1}^{N_k}\log p_j^{(k)}\,|f(p_j^{(k)})|^2.
\]
\end{definition}

\begin{lemma}[Weighted trace identity]\label{lem:weighted}
For every $k$,
\[
  \operatorname{tr}(\GVMkw)
  =\sum_{j=1}^{N_k}\log p_j^{(k)}\,K_{N_k}(p_j^{(k)},p_j^{(k)}).
\]
Moreover, if Assumption~\textup{(H)} holds and the weighted analogue of the remainder
estimate is available along the sequence, then the weighted form admits a comparison
with the reference operator up to an error term controlled by $\log p_{N_k}\,\varepsilon_{N_k}$.
\end{lemma}

\begin{proof}
The trace identity is immediate from the definition of $\GVMkw$.
The second statement is formal: it is obtained by inserting the logarithmic weights into
the same matrix comparison scheme as in Paper~I, provided the corresponding weighted
error control is assumed.
We do not claim this weighted comparison unconditionally.
\end{proof}

\subsection{Asymptotic trace formulae}

\begin{theorem}[Trace formula]\label{thm:trace}
Let $(N_k,c_k,\omega_k,T_k,\mathbf p_k)$ be a $\rho$-admissible prime sequence.
Assume that the diagonal kernel asymptotic is valid along the sampled prime points in the sense that
\[
  K_{N_k}(p_j^{(k)},p_j^{(k)})
  =\frac{\omega_k N_k}{\pi}\bigl(1+r_{k,j}\bigr)
\]
with an averaged error term satisfying
\[
  \frac1{N_k}\sum_{j=1}^{N_k}|r_{k,j}|\to0.
\]
Then, as $k\to\infty$:
\begin{enumerate}[(I)]
\item \textbf{Unweighted trace:}
\[
  \frac{\pi}{\omega_kN_k^2}\operatorname{tr}(\GVMk)\to1.
\]
\item \textbf{Weighted trace:}
\[
  \frac{\pi}{\omega_kN_k^2\log N_k}\operatorname{tr}(\GVMkw)\to1.
\]
\item \textbf{Prime-number-theorem consistency:}
\[
  \frac{1}{N_k\log N_k}
  \sum_{j=1}^{N_k}\log p_j^{(k)}\cdot \frac{\pi}{\omega_k}K_{N_k}(p_j^{(k)},p_j^{(k)})
  \to1.
\]
\end{enumerate}
\end{theorem}

\begin{proof}
For part (I), Lemma~\ref{lem:dual_trace}(i) and the averaged kernel asymptotic give
\[
  \operatorname{tr}(\GVMk)
  =\sum_{j=1}^{N_k}\frac{\omega_kN_k}{\pi}(1+r_{k,j})
  =\frac{\omega_kN_k^2}{\pi}\Bigl(1+\frac1{N_k}\sum_{j=1}^{N_k}r_{k,j}\Bigr),
\]
which implies the limit.

For part (II), we similarly obtain
\[
  \operatorname{tr}(\GVMkw)
  =\frac{\omega_kN_k}{\pi}
  \sum_{j=1}^{N_k}\log p_j^{(k)}\,(1+r_{k,j}).
\]
By the prime number theorem in the form
$\sum_{j=1}^{N_k}\log p_j^{(k)}\sim N_k\log N_k$, the main term is
$\omega_kN_k^2\log N_k/\pi$.
The averaged error remains negligible under the stated hypothesis.
This proves (II).

Part (III) is merely a rewriting of (II).
\end{proof}

\begin{remark}[Status of the trace input]
Theorem~\ref{thm:trace} is intentionally formulated so that the analytic input is explicit.
The arithmetic input is the prime number theorem, while the PSWF input is an averaged diagonal
kernel asymptotic along the prime sample points.
This avoids hiding an additional unproved statement inside the normalization.
\end{remark}

\begin{remark}[Relation to the Weil form]\label{rem:weil_compare}
The weighted trace is only a trace-level analog of the prime part of Weil's explicit formula.
To pass from this averaged operator statement to an actual Weil operator, one would need a
Mellin-side identification of the relevant test-function geometry.
That structural step is not supplied here.
\end{remark}

\section{Open Problems}\label{sec:open}

\begin{problem}[Prime sampling and DSTP]\label{prob:primeT2}
Prove or disprove Assumption~\textup{(H)} for prime sampling.
More precisely, determine whether the prime sets
$\mathbf p_k=(2,3,5,\ldots,p_{N_k})$ satisfy the DSTP of \cite{PaperI}
with defect level $\varepsilon_{N_k}\to0$ in the regime $N_k/c_k\to\rho\in(0,2/\pi)$.
Any unconditional coercivity theorem for prime sampling in the present scaling limit depends on this.
\end{problem}

\begin{problem}[Trace preservation and uniqueness]\label{prob:unique}
Show that the normalized operators $\widetilde{\mathcal G}_k$ have a unique weak-$*$ limit and that
$\operatorname{tr}(\Gcal_\infty)=1$.
This requires ruling out escape of trace mass at infinity.
\end{problem}

\begin{problem}[Mellin-isometric embedding]\label{prob:mellin}
Construct an isometric embedding from the additive $L^2(\mathbb R)$ framework of the PSWFs
into the multiplicative space relevant to Weil's explicit formula, in such a way that
$\Gcal_\infty$ corresponds to a genuine Weil operator.
\end{problem}

\begin{problem}[Weil positivity and RH]\label{prob:positivity}
Assuming Problems~\ref{prob:unique} and \ref{prob:mellin}, determine whether the positivity of the
resulting limit operator is equivalent to the Riemann Hypothesis via Weil's criterion.
\end{problem}

\begin{problem}[Critical density $\rho=2/\pi$]\label{prob:critical}
Analyze the critical regime $N_k/c_k\to2/\pi$, where the Landau--Widom bulk estimate degenerates.
\end{problem}

\section*{Acknowledgements}
The author thanks [---] for helpful discussions.
Paper~I of this program is \cite{PaperI}.

\begin{thebibliography}{99}

\bibitem{PaperI}
Anonymous Author,
\emph{Coercivity of the Prolate Gram Form at Discrete Sampling Points},
preprint, 2026.

\bibitem{Bombieri2000}
E.~Bombieri,
\textit{Problems of the Millennium: The Riemann Hypothesis},
Clay Mathematics Institute, 2000.

\bibitem{BonamiKaroui}
A.~Bonami and A.~Karoui,
\textit{Spectral decay of the sinc kernel operator and approximation
with prolate spheroidal wave functions},
arXiv:1012.3881.

\bibitem{CC2019}
A.~Connes and C.~Consani,
\textit{Riemann--Roch and Weil positivity},
arXiv:1910.14368 (2019).

\bibitem{CCM_PNAS2022}
A.~Connes, C.~Consani, and H.~Moscovici,
\textit{The UV prolate spectrum matches the zeros of zeta},
Proc.\ Natl.\ Acad.\ Sci.\ \textbf{119} (2022), e2123174119.

\bibitem{CCM2025}
A.~Connes, C.~Consani, and H.~Moscovici,
arXiv:2511.22755 (2025).

\bibitem{ReedSimon}
M.~Reed and B.~Simon,
\textit{Methods of Modern Mathematical Physics~I},
Academic Press, 1980.

\bibitem{Weil1952}
A.~Weil,
\textit{Sur les formules explicites de la th\'{e}orie des nombres},
Izv.\ Akad.\ Nauk \textbf{36} (1952), 3--18.

\end{thebibliography}

\end{document}
